% Generated draft: requires mathematical and source review.
\documentclass[11pt]{article}
\usepackage[margin=25mm]{geometry}
\usepackage{fontspec}
\setmainfont{Latin Modern Roman}
\usepackage{amsmath,amssymb,mathtools,unicode-math}
\setmathfont{Latin Modern Math}
\usepackage{enumitem,needspace,xurl,hyperref}
\setlength{\parindent}{0pt}
\setlength{\parskip}{5pt}
\setlength{\emergencystretch}{3em}
\newcommand{\N}{\mathbb{N}}
\newcommand{\Z}{\mathbb{Z}}
\newcommand{\Q}{\mathbb{Q}}
\newcommand{\R}{\mathbb{R}}
\newcommand{\C}{\mathbb{C}}
\newcommand{\F}{\mathbb{F}}
\newcommand{\A}{\mathbb{A}}
\newcommand{\PP}{\mathbb{P}}
\newcommand{\E}{\mathbb{E}}
\newcommand{\eps}{\varepsilon}
\newcommand{\dd}{\,\mathrm{d}}
\begin{document}
\section*{UnsolvedMath ID: 30002390}
\textbf{Source problem code:} catalogue entry
\subsection{Definitions and mathematical statement}
Let \(\mathbb H^n\) denote real hyperbolic \(n\)-space of sectional curvature \(-1\), and identify its ideal boundary with \(S^{n-1}\). For a subgroup \(\Gamma<\operatorname{Isom}(\mathbb H^4)\), define its limit set by
\[
\Lambda(\Gamma)=\overline{\Gamma x}\cap\partial_\infty\mathbb H^4,
\qquad x\in\mathbb H^4,
\]
where the closure is taken in the ideal compactification. Write \(\dim_H\Lambda(\Gamma)\) for its Hausdorff dimension in the standard round metric on \(\partial_\infty\mathbb H^4\cong S^3\). A hyperbolic manifold is a Riemannian manifold of constant sectional curvature \(-1\). The source does not specify completeness, compactness, or volume restrictions for the manifolds below.

\textbf{Limit-set question.} Interpreting a hyperbolic \(3\)-manifold group as a group isomorphic to the fundamental group of a connected hyperbolic \(3\)-manifold, and provisionally interpreting the displayed action as discrete, is the following assertion true?
\[
\begin{gathered}
\text{For every }\varepsilon>0\text{, there exist a connected hyperbolic }3\text{-manifold }N
\text{ and a discrete subgroup }\Gamma<\operatorname{Isom}(\mathbb H^4)\\
\text{such that }\Gamma\cong\pi_1(N)
\quad\text{and}\quad
3-\varepsilon<\dim_H\Lambda(\Gamma)\leq 3.
\end{gathered}
\]
The subgroup may depend on \(\varepsilon\). Discreteness and the intended geometric class of \(N\) require confirmation from the original problem context. This formulation allows dimension exactly \(3\); the source does not say whether dimensions strictly below \(3\) are required.

\textbf{Fibration question.} Does there exist a connected real hyperbolic \(4\)-manifold \(M\), a connected surface \(B\), and a surface \(F\), together with a locally trivial fiber bundle
\[
F\hookrightarrow M\xrightarrow{\ p\ }B?
\]
Here interpreting ``fibering'' as a locally trivial fiber bundle is provisional: the source does not specify the smooth or topological category. It also does not specify completeness or finite volume of \(M\), compactness, orientability, genus, or boundary conditions for the manifolds and surfaces. These missing qualifications remain unresolved and may affect the intended existence question.

The source asks about the full limit set \(\Lambda(\Gamma)\), whereas its supplied literature assessment refers to a conical limit set. These are distinct notions, so that assessment does not justify replacing the limit set in the question. The two existence questions are retained separately; no equivalence or implication between them is asserted.
\subsection{Short English statement}
The first question asks whether hyperbolic \(3\)-manifold groups acting on \(\mathbb H^4\) can have limit-set Hausdorff dimensions arbitrarily close to \(3\). The second asks whether a real hyperbolic \(4\)-manifold can be a surface bundle over a surface. The intended geometric hypotheses and bundle category need clarification. Both questions are labeled open in the supplied record; this draft does not independently verify that status.
\subsection{Sources}
\begin{itemize}
\item Local UnsolvedMath registry, permanent ID \(30002390\), source problem code \texttt{catalogue entry}, title ``Hyperbolic Limit Sets and Fibered Four-Manifolds,'' in the Oberwolfach Reports Open Problems collection. No individual proposer is supplied; the proposed year is recorded as 2013.
\item Supplied report citation: \emph{Group Theory, Measure, and Asymptotic Invariants} (2014), Oberwolfach Reports, pp.~2375--2422, DOI \texttt{10.4171/owr/2013/42}. \url{https://doi.org/10.4171/owr/2013/42}. The supplied background describes the report as documenting workshop discussions and open problems. Its citation year is 2014, while the proposed year and DOI contain 2013; this draft preserves those supplied details without resolving them.
\item Supplied literature reference, labeled ``Bowen primary paper'': \url{https://arxiv.org/abs/1303.5963}. No paper title or precise theorem reference is supplied. The literature assessment dated \texttt{2026-08-21} states that Bowen proves a uniform upper bound for the relevant conical limit-set dimension and labels the near-\(3\) construction and the related fibration problem open. This is supplied dated curation, not independently verified literature assessment; the record itself requests independent verification of time-sensitive status and cited claims.
\end{itemize}
\end{document}
